% Source: https://artofproblemsolving.com/community/c3h1366899_geometry_medians_problem
% Stable identity: aops-topic-1366899-q01
Let the medians meet at the centroid $G$, which divides each median in the ratio $2:1$ from the vertex. So $BG=4$, $GD=2$, $CG=6$, $GE=3$.

\begin{asy}
usepackage("amsmath");
size(5.5cm);
defaultpen(fontsize(10pt));
path ra(pair P,pair u,pair v,real s=0.2){return P+s*unit(u)--P+s*unit(u)+s*unit(v)--P+s*unit(v);}
pair ft(pair P,pair A,pair B){return A+dot(P-A,unit(B-A))*unit(B-A);}
pair G=(0,0),B=(-4,0),D=(2,0),C=(0,-6),EE=(0,3),A=(4,6);
draw(A--B--C--cycle,linewidth(0.9));
draw(B--D,red); draw(C--EE,blue);
draw(ra(G,D-G,EE-G,0.5));
dot(G); dot(D); dot(EE);
label("$A$",A,NE); label("$B$",B,W); label("$C$",C,S); label("$D$",D,E); label("$E$",EE,NW); label("$G$",G,SW);
label("$4$",(B+G)/2,S,red); label("$2$",(G+D)/2,S,red); label("$6$",(C+G)/2,E,blue); label("$3$",(G+EE)/2+(0,0.3),E,blue);
label("\mbox{}",(min(currentpicture,user=true)+max(currentpicture,user=true))/2,black); // keeps the text after the figure black when LaTeX typesets the labels inline
\end{asy}

Since $BD\perp CE$, triangle $BGC$ has a right angle at $G$ and area $\dfrac{1}{2}\cdot4\cdot6=12$.

The three medians cut a triangle into six small triangles of equal area, so $[BGC]$ (two of them) is $\dfrac{1}{3}$ of $[ABC]$. Hence
\[[ABC]=3\cdot12=\boxed{36}.\]
(Coordinate check: with $G$ at the origin, $B=(-4,0)$, $D=(2,0)$, $C=(0,-6)$, $E=(0,3)$, we get $A=2D-C=(4,6)$, and indeed $E$ is the midpoint of $AB$; the shoelace formula gives area $36$.) A useful general fact: if two medians of lengths $m$ and $n$ are perpendicular, the area is $\dfrac{2}{3}mn$.
